\section*{Chapitre \ref{ch:g12:internal-heat} --- La chaleur interne de la Terre}

\begin{solution}{exo:g12:internal-heat:1}
(a) le \omterm{def:g12:internal-heat:heat-flow}{flux géothermique}~; (b) la \omterm{def:g12:internal-heat:heat-production}{production de chaleur radiogénique}~;
(c) le \omterm{def:g12:internal-heat:secular-cooling}{refroidissement séculaire}~; (d) la \omterm{def:g12:internal-heat:convection}{convection thermique}~;
(e) l'\omterm{def:g12:internal-heat:geothermal}{énergie géothermique}.
\end{solution}

\begin{solution}{exo:g12:internal-heat:2}
(1) \qty{0.0916}{W/m^2}. (2) \qty{4.67e13}{W}. (3) \qty{0.029}{K/m}~: une
différence d'un degré Celsius est une différence d'un kelvin.
\end{solution}

\begin{solution}{exo:g12:internal-heat:3}
Le thorium~232, l'uranium~238, le potassium~40 et l'uranium~235. KamLAND
ne voit que les géoneutrinos du thorium~232 et de l'uranium~238.
\end{solution}

\begin{solution}{exo:g12:internal-heat:4}
$490/\sqrt{1} = \qty{490}{mW/m^2}$~; $490/2 = \qty{245}{mW/m^2}$~;
$490/5 = \qty{98}{mW/m^2}$~; $490/7 = \qty{70}{mW/m^2}$, à peu près la
moyenne des continents.
\end{solution}

\begin{solution}{exo:g12:internal-heat:5}
Près de la surface, la température du sol suit les jours et les saisons~;
ce n'est qu'au-dessous de quelques mètres qu'elle reste à la température
moyenne de l'année, et plus profond encore qu'elle révèle la chaleur venue
d'en bas. Le forage brasse la roche et l'eau du trou et modifie leur
température~: on attend donc que le trou soit revenu à la température de la
roche.
\end{solution}

\begin{solution}{exo:g12:internal-heat:6}
(1) $(265 - 10)/9.1 \approx \qty{28.0}{\celsius}$ par kilomètre, soit
\qty{0.0280}{K/m}. (2) $q = 2.72 \times 0.0280 \approx \qty{0.0762}{W/m^2}$,
environ \qty{76}{mW/m^2}. (3) Un peu au-dessus de la moyenne des
continents, \qty{70.9}{mW/m^2}, d'environ \qty{7}{\percent}~: une valeur
continentale typique.
\end{solution}

\begin{solution}{exo:g12:internal-heat:7}
Continents~: $0.0709 \times \num{2.073e14} \approx \qty{1.47e13}{W} =
\qty{14.7}{TW}$. Océans~: $0.1054 \times \num{3.028e14} \approx
\qty{3.19e13}{W} = \qty{31.9}{TW}$. Total~: environ \qty{46.6}{TW}, dont
$31.9/46.6 \approx \qty{68}{\percent}$, soit les deux tiers, sortent par
les fonds océaniques.
\end{solution}

\begin{solution}{exo:g12:internal-heat:8}
(1) $\num{1.88e-6} \times 8000 \approx \qty{0.0150}{W/m^2} = \qty{15.0}{mW/m^2}$.
(2) $15.0/70.9 \approx \qty{21}{\percent}$, le cinquième du \omterm{def:g12:internal-heat:heat-flow}{flux
géothermique}.
\end{solution}

\begin{solution}{exo:g12:internal-heat:9}
$\sqrt{t} = 490/70.9 \approx 6.91$, donc $t \approx 47.8$ millions
d'années~: les fonds océaniques âgés de plus de 48 millions d'années
environ perdent moins de chaleur que le continent moyen.
\end{solution}

\begin{solution}{exo:g12:internal-heat:10}
Uranium~235~: $2^{1.4/0.704} = 2^{1.99} \approx 4$ fois plus.
Potassium~40~: $2^{2.5/1.248} = 2^{2.00} \approx 4$ fois plus. Dans les
deux cas, on remonte de deux demi-vies.
\end{solution}

\begin{solution}{exo:g12:internal-heat:11}
(1) $10^2/10^{-6} = \qty{1e8}{s}$, environ 3 ans. (2) $(10^5)^2/10^{-6} =
\qty{1e16}{s}$, environ 320 millions d'années. (3) $(\num{2.891e6})^2/10^{-6}
\approx \qty{8.4e18}{s}$, environ 265 milliards d'années. Le temps croît
comme le carré de la distance~: dix fois plus épais, cent fois plus long.
\end{solution}

\begin{solution}{exo:g12:internal-heat:12}
Multiplier chacun par $2^{3/T}$~: thorium~232, $10.4 \times 2^{3/14.0} \approx
\qty{12.1}{TW}$~; uranium~238, $9.6 \times 2^{3/4.468} \approx \qty{15.3}{TW}$~;
potassium~40, $4.0 \times 2^{3/1.248} \approx \qty{21.2}{TW}$~; uranium~235,
$0.42 \times 2^{3/0.704} \approx \qty{8.1}{TW}$. Total~: environ
\qty{56.7}{TW}, plus du double des \qty{24.4}{TW} actuels. Les deux
isotopes à vie courte en fournissent
$(21.2 + 8.1)/56.7 \approx \qty{52}{\percent}$, contre
$4.42/24.4 \approx \qty{18}{\percent}$ aujourd'hui.
\end{solution}

\begin{solution}{exo:g12:internal-heat:13}
La puissance est le produit du débit massique, de la capacité thermique
massique et du réchauffement~: $\num{3e8} = 1640 \times 4180 \times \Delta T$, donc $\Delta T =
\num{3e8}/(1640 \times 4180) \approx \qty{44}{K}$. Le continent moyen perd
\qty{0.0709}{W} par mètre carré~; \qty{300}{MW} représentent donc la perte
naturelle de chaleur de $\num{3e8}/0.0709 \approx \qty{4.2e9}{m^2}$,
environ \qty{4200}{km^2}~: la centrale recueille une chaleur que la Terre
laisse échapper naturellement sur des milliers de kilomètres carrés.
\end{solution}

\begin{solution}{exo:g12:internal-heat:14}
(1) 100 millions d'années, c'est $\num{2.6e11}/\num{1e8} = 2600$ fois
moins que 260 milliards d'années~; 200 millions d'années, 1300 fois moins.
(2) La convection transporte la chaleur avec la roche en mouvement~: le
temps qu'il faut pour traverser le manteau est son épaisseur divisée par la
vitesse, $\qty{2891}{km}$ à quelques centimètres par an, soit une centaine
de millions d'années. La conduction n'a pas de tel transporteur~: son temps
croît comme le carré de la distance, et sur des milliers de kilomètres il
devient énorme.
\end{solution}

\begin{solution}{exo:g12:internal-heat:15}
(1) $10/22.3 \approx \qty{45}{\percent}$~; le reste est le refroidissement
du manteau lui-même, qui perd sa propre chaleur emmagasinée. (2) La limite
se trouve au rayon $6371 - 2891 =
\qty{3480}{km}$~; sa surface vaut $4\pi (\num{3.48e6})^2 \approx
\qty{1.52e14}{m^2}$, et $\num{1e13}/\num{1.52e14} \approx
\qty{0.066}{W/m^2}$, environ \qty{66}{mW/m^2}, quelque \qty{72}{\percent}
de la moyenne en surface. (3) La chaleur traverse la limite par conduction
jusque dans la couche la plus basse du manteau~; réchauffée, cette roche
devient plus légère que le manteau situé au-dessus et peut monter, en
emportant la chaleur vers le haut~: c'est ainsi que naîtraient les
panaches, en profondeur.
\end{solution}

\begin{solution}{pb:g12:internal-heat:1}
\textbf{1.} $4\pi (\num{6.371e6})^2 \approx \qty{5.101e14}{m^2}$, et
$2.073 + 3.028 = 5.101$~: oui.

\textbf{2.} $3.028/5.101 \approx \qty{59}{\percent}$, environ les trois
cinquièmes.

\textbf{3.} Les fonds océaniques sont jeunes~: ils naissent chauds aux
dorsales et se refroidissent en s'éloignant, et une \omterm{def:g11:lithosphere-asthenosphere:lithosphere}{lithosphère} jeune perd
beaucoup de chaleur.

\textbf{4.} $490/\sqrt{4} = \qty{245}{mW/m^2}$.

\textbf{5.} $2.72 \times 0.026 \approx \qty{0.0707}{W/m^2} =
\qty{70.7}{mW/m^2}$~: juste en dessous de la moyenne, pratiquement égal à
elle.

\textbf{6.} $0.0709 \times \num{2.073e14} \approx \qty{14.7}{TW}$.

\textbf{7.} $0.1054 \times \num{3.028e14} \approx \qty{31.9}{TW}$.

\textbf{8.} $14.7 + 31.9 \approx \qty{46.6}{TW}$, les \qty{47}{TW} de
l'estimation publiée (la moyenne mondiale du chapitre donne
\qty{46.7}{TW}, la valeur utilisée dans la suite).

\textbf{9.} $\num{4.66e13}/\num{5.101e14} \approx \qty{0.0914}{W/m^2}$,
soit environ \qty{91}{mW/m^2}.

\textbf{10.} $0.0243 \times \num{2.073e14} \approx \qty{5.0}{TW}$, le tiers
des \qty{14.7}{TW} perdus à travers les continents~: la radioactivité de la
\omterm{def:g12:continental-crust:continental-crust}{croûte continentale} fournit une grande part de leur \omterm{def:g12:internal-heat:heat-flow}{flux géothermique}.

\textbf{11.} Pour chaque kilogramme d'uranium, la Terre contient
\qty{3.9}{kg} de thorium, qui libèrent $3.9 \times 26.3 \approx \qty{102.6}{\micro W}$,
contre \qty{94.4}{\micro W} pour l'uranium~238. L'uranium~238 produit donc
$94.4/(94.4 + 102.6) \approx 0.479$ des \qty{20.0}{TW}~: environ
\qty{9.6}{TW}. Uranium~235~: $9.59 \times 4.10/94.4 \approx \qty{0.42}{TW}$.

\textbf{12.} $20.0 + 0.42 + 4 \approx \qty{24.4}{TW}$.

\textbf{13.} Le granite et les autres roches de la \omterm{def:g12:continental-crust:continental-crust}{croûte continentale},
environ vingt fois plus riches que les roches ultramafiques du manteau.

\textbf{14.} Plus grande~: chaque isotope était alors plus abondant, deux
fois plus pour chaque demi-vie en arrière~; il y a 3~milliards d'années,
le total valait plus du double de la valeur actuelle.

\textbf{15.} $46.7 - 24.4 = \qty{22.3}{TW}$.

\textbf{16.} De la chaleur emmagasinée depuis la formation de la Terre~:
c'est le \omterm{def:g12:internal-heat:secular-cooling}{refroidissement séculaire} du manteau et du noyau.

\textbf{17.} $22.3/46.7 \approx \qty{48}{\percent}$.

\textbf{18.} De $11/46.7 \approx \qty{24}{\percent}$ à $38/46.7 \approx
\qty{81}{\percent}$~: les géoneutrinos resserrent beaucoup cet intervalle.

\textbf{19.} La \omterm{def:g12:internal-heat:convection}{convection mantellique}, qui transporte l'essentiel de la
chaleur~; seules les dernières dizaines de kilomètres sont traversées par
conduction.

\textbf{20.} $24.4/46.7 \approx 0.52$~: la radioactivité fournit
aujourd'hui \textbf{environ la moitié} de la perte de chaleur de la Terre,
et le lent refroidissement de la planète l'autre moitié.
\end{solution}
